\section{拍频与近频率叠加}

同方向、同振幅而角频率接近的两列简谐波在固定位置可写成
\begin{align}
 y&=A\cos(\omega_1t+\phi_1)+A\cos(\omega_2t+\phi_2)\\
 &=2A\cos\!\left(\frac{\Delta\omega}{2}t+\frac{\Delta\phi}{2}\right)
 \cos\!\left(\bar\omega t+\bar\phi\right),
\end{align}
其中 $\Delta\omega=\omega_1-\omega_2$、$\bar\omega=(\omega_1+\omega_2)/2$。
快速振动的包络振幅缓慢变化。相邻两次响度极大的时间间隔对应包络绝对值的周期，故拍频
\begin{equation}
 f_b=|f_1-f_2|.
\end{equation}
拍频不是第三个独立本征频率，而是线性叠加后能量或响度包络的变化频率。

\begin{center}
\begin{tikzpicture}
\begin{axis}[width=.82\linewidth,height=5.0cm,axis lines=middle,
  xlabel={$t$},ylabel={$y$},ylabel style={at={(axis description cs:-.025,.5)},anchor=east},xmin=0,xmax=12.6,ymin=-2.3,ymax=3.2,ticks=none,
  every axis label/.style={fill=none}]
  \addplot[DeepBlue,very thick,domain=0:12.6,samples=650]{2*cos(deg(.5*x))*cos(deg(7*x))};
  \addplot[TextGray,dashed,domain=0:12.6,samples=180]{2*abs(cos(deg(.5*x)))};
  \addplot[TextGray,dashed,domain=0:12.6,samples=180]{-2*abs(cos(deg(.5*x)))};
  \draw[<->] (axis cs:0,2.55)--(axis cs:6.283185,2.55) node[midway,above] {$T_b=1/|f_1-f_2|$};
\end{axis}
\end{tikzpicture}
\figcap{两近频率简谐振动叠加形成拍。虚线是振幅包络，不是介质中另一列波。}\label{ch12:fig:beats}
\end{center}
图~\ref{ch12:fig:beats} 中相邻两个包络极大之间就是一个拍周期，快速载波在此期间振动许多次。

\begin{origexamplebox}{用拍频校准音叉}
\label{ch12:ex:03}
标准音叉频率为 $440\,\mathrm{Hz}$，待测音叉与它同时发声时每秒听到3次响度增强。在待测音叉上粘少量蜡后，拍频降为 $1\,\mathrm{Hz}$。求待测音叉原频率。
\end{origexamplebox}
\begin{xesolution}
初始只有两个候选：
\[
 f=440\pm3\,\mathrm{Hz},
\]
即 $437\,\mathrm{Hz}$ 或 $443\,\mathrm{Hz}$。粘蜡增加振动系统的有效质量，使待测频率降低。若原来是
$443\,\mathrm{Hz}$，降低后会接近 $440\,\mathrm{Hz}$，拍频可从3降到1；若原来是 $437\,\mathrm{Hz}$，继续降低只会使频差增大。
因此
\[
 \boxed{f=443\,\mathrm{Hz}}.
\]
方向判断是这道题不可省略的一步；单凭初始拍频不能确定高低。
\end{xesolution}

\section{多普勒效应}

声波相对介质传播。这里讨论均匀静止介质中的共线运动、声源亚声速 $|v_s|<c$，所有速度均相对介质测量。设介质中声速为 $c$，观察者以速度 $v_o$ 向声源运动，声源以速度 $v_s$ 向观察者运动，则
\begin{equation}
 f'=f\frac{c+v_o}{c-v_s}.
\end{equation}
若实际运动方向相反，应给相应速度取负号。分子来自观察者遇到波峰的速率变化，分母来自移动声源改变波长。
对于真空光，不能套用“相对介质速度”的声学公式，应使用相对论多普勒关系。

\begin{center}
\begin{tikzpicture}[scale=.85,>=Latex]
  \foreach \n in {1,2,3,4}{
   \draw[TextGray] ({-.245*\n},0) circle({.7*\n});
   \fill[TextGray] ({-.245*\n},0) circle(1pt);
  }
  \fill[DeepBlue] (0,0) circle(2.8pt);
  \node[below] at (0,-.12) {$S$};
  \fill (3.9,0) circle(2.8pt) node[right=4pt] {$O$};
  \draw[vector] (.1,-3.3)--(1.3,-3.3) node[midway,below] {$v_s$};
  \draw[vector] (3.8,-1.1)--(2.6,-1.1) node[midway,below] {$v_o$};
  \draw[vector] (.3,3.1)--(1.7,3.1) node[midway,above] {$c$};
  \draw[<->] (1.365,-.30)--(1.82,-.30);
  \draw[guide] (1.60,-.35)--(2.6,-2.35);
  \node[right] at (2.7,-2.45) {$\lambda_{\text{前}}=(c-v_s)/f$};
  \draw[<->] (-3.78,-.35)--(-2.835,-.35);
  \node[below] at (-3.3,-.45) {$\lambda_{\text{后}}$};
  \node at (-1.5,-4.3) {圆心是历次发射位置，非当前声源位置};
\end{tikzpicture}
\figcap{声源和观察者相向运动的符号约定。波面相对介质以声速 $c$ 传播。}\label{ch12:fig:doppler}
\end{center}
图~\ref{ch12:fig:doppler}中，第 $n$ 个波面发出于 $-n/f$ 时，
圆心为 $x=-nv_s/f$、半径为 $nc/f$。相邻前方波峰间距为
$(c-v_s)/f$，观察者迎着波峰运动，单位时间跨过 $(c+v_o)/\lambda_{\text{前}}$ 个波峰，
这就推出上面的频率公式。图固定了正号方向：$v_s,v_o>0$ 都使接收频率升高。

\begin{origexamplebox}{迎面驶来的鸣笛列车}
\label{ch12:ex:04}
列车以 $20\,\mathrm{m/s}$ 接近静止观察者，汽笛本征频率 $600\,\mathrm{Hz}$，空气中声速
$340\,\mathrm{m/s}$。求接近时和驶离后的听到频率，忽略风。
\end{origexamplebox}
\begin{xesolution}
观察者静止，$v_o=0$。接近时声源向观察者，
\[
 f_{\rm app}=600\frac{340}{340-20}
 =\boxed{637.5\,\mathrm{Hz}}.
\]
驶离时 $v_s$ 取负，
\[
 f_{\rm rec}=600\frac{340}{340+20}
 =\boxed{566.7\,\mathrm{Hz}}.
\]
接近频率高、驶离频率低，符合波面前方被压缩、后方被拉长的图像。分式分子、分母都是速度，故频率单位保持不变。
\end{xesolution}

\section{声强与声强级}

声强 $I$ 是单位面积上平均能流功率，单位为 $\mathrm{W/m^2}$。点声源在均匀无耗介质中各向同性辐射时
\begin{equation}
 I(r)=\frac{P}{4\pi r^2}.
\end{equation}
声强级定义为
\begin{equation}
 L_I=10\log_{10}\frac I{I_0}\quad(\mathrm{dB}),
 \qquad I_0=10^{-12}\,\mathrm{W/m^2}.
\end{equation}
分贝是对数比值；多个互不相干声源应先相加强度，再取对数，不能直接相加分贝数。

\begin{origexamplebox}{相同声源的声强级叠加}
\label{ch12:ex:05}
一台机器在某点产生 $70\,\mathrm{dB}$ 的声强级。若有8台相同机器独立运行，且它们在该点的声波相位互不相关，求总声强级。
\end{origexamplebox}
\begin{xesolution}
每台强度相同，非相干时总强度为 $I_{\rm tot}=8I$。因此
\[
 L_{\rm tot}-L_1=10\log_{10}8\approx9.03\,\mathrm{dB}.
\]
故
\[
 \boxed{L_{\rm tot}\approx79.0\,\mathrm{dB}}.
\]
不是 $8\times70\,\mathrm{dB}$。若只增加到两台，相同推理给出约 $3.01\,\mathrm{dB}$ 的增量。
\end{xesolution}

\subsection{声强与声压、质元速度怎样联系}
声音传播时，空气质元在平衡位置附近往复运动，并不是从声源一路飞到听者。取沿 $+x$ 传播的平面简谐声波，位移为
\[
 \xi(x,t)=\xi_m\cos(kx-\omega t),\qquad \omega=ck.
\]
质元速度是对时间求偏导，而波速是波形的传播速度，二者不能混淆：
\[
 v_{\rm part}=\partial_t\xi=\omega\xi_m\sin(kx-\omega t).
\]
小形变的体压缩量为 $-\partial_x\xi$；用绝热体积模量 $B_s=\rho_0c^2$，声压扰动为
\[
 p'=-B_s\partial_x\xi
 =\rho_0c\omega\xi_m\sin(kx-\omega t)
 =\rho_0c\,v_{\rm part}.
\]
这里 $p'$ 是总压强相对平衡值的微小变化，可以为正或负；$\rho_0$ 为平衡密度。对于这列右行波，声压与质元速度同相。截面上每单位面积传递的瞬时功率为 $p'v_{\rm part}$，取一个周期的平均，得到
\[
 I=\langle p'v_{\rm part}\rangle
 =\frac12\rho_0c\omega^2\xi_m^2
 =\frac{p_m^2}{2\rho_0c}
 =\frac{p_{\rm rms}^2}{\rho_0c}.
\]
因正弦平方的周期平均为 $1/2$，有效声压 $p_{\rm rms}=p_m/\sqrt2$。所以声强与振幅的平方成正比，而不是与振幅本身成正比。上述声压与声强关系适用于单向平面行波；在驻波中声压和速度的时空关系不同，不能照搬。

\subsection{先算强度，再换成分贝}
球面面积随距离平方增加。在无吸收且各向同性辐射的条件下，每个球面通过的总功率都为 $P$，所以 $4\pi r^2I=P$。这不是能量丢失，而是同一功率分散到更大的面积。距离从 $r_1$ 变成 $r_2$ 时，
\[
 L_I(r_2)-L_I(r_1)
 =10\log_{10}\frac{I(r_2)}{I(r_1)}
 =-20\log_{10}\frac{r_2}{r_1}.
\]
距离加倍，声强减为四分之一，声强级下降约 $6.02\,\mathrm{dB}$。这与“两个独立相同声源使声强级增加约 $3.01\,\mathrm{dB}$”并不矛盾，前者改变传播面积，后者改变总入射强度。

\begin{origexamplebox}{不同距离、不同声强级的两个独立声源}\label{ch12:ex:teaching-sound-distance}
两个可视为点声源的机器各向同性发声。在距机器甲 $2\,\mathrm m$ 处，它单独产生 $80\,\mathrm{dB}$ 的声强级；机器乙在观察点单独产生 $70\,\mathrm{dB}$。观察点距甲 $4\,\mathrm m$。忽略反射和吸收，二者互不相干。求同时工作时观察点的总声强级。
\end{origexamplebox}
\begin{xesolution}
先把甲的测量位置换到观察点。距离加倍，强度减为四分之一，故
\[
 I_A=\frac14I_0 10^{80/10}=2.5\times10^{-5}\,\mathrm{W/m^2}.
\]
乙在该点的强度为 $I_B=I_0 10^{70/10}=10^{-5}\,\mathrm{W/m^2}$。互不相干时，平均交叉项为零，才有 $I=I_A+I_B$，所以
\[
 L=10\log_{10}(3.5\times10^7)\approx75.44\,\mathrm{dB}.
\]
结果略高于甲单独在该点的 $73.98\,\mathrm{dB}$，不是两个分贝数的算术和。若乙关闭，保留 $I_A$ 即回到甲单独的结果。
\end{xesolution}

若两列同频同向平面声波保持固定相位差 $\delta$，应先叠加声压，再平方取平均，得到
\[
 I=I_1+I_2+2\sqrt{I_1I_2}\cos\delta.
\]
等强同相时局部强度为 $4I_1$，反相时可以相消；若相位不断随机变化，平均 $\cos\delta=0$，才恢复非相干相加。驻波、不同传播方向和复杂声场还应计算矢量能流，不能把这条同向行波公式当作任意声场的通式。

\section{波包、相速度与群速度}

把相近波数的简谐波叠加可形成局域波包。单个相位点按
\begin{equation}
 v_p=\frac\omega k
\end{equation}
传播，而窄带包络按
\begin{equation}
 v_g=\frac{\dd\omega}{\dd k}
\end{equation}
传播。无色散介质中 $\omega=ck$，二者相等；有色散时二者通常不同。若波包带宽太大或介质强吸收，单一群速度未必能完整描述包络变形。

\begin{center}
\begin{tikzpicture}
\begin{axis}[width=.82\linewidth,height=4.9cm,axis lines=middle,
  xlabel={$x$},ylabel={$y$},ylabel style={at={(axis description cs:-.025,.5)},anchor=east},xmin=-6,xmax=6,ymin=-1.5,ymax=1.7,ticks=none,
  every axis label/.style={fill=none}]
  \addplot[DeepBlue,very thick,domain=-6:6,samples=650]{exp(-x^2/8)*cos(deg(5*x))};
  \addplot[TextGray,dashed,domain=-6:6,samples=220]{exp(-x^2/8)};
  \addplot[TextGray,dashed,domain=-6:6,samples=220]{-exp(-x^2/8)};
  \draw[vector] (axis cs:1.2,1.35)--(axis cs:3.5,1.35) node[midway,above,clearlabel]{$v_g$};
  \draw[component] (axis cs:2.0,-.85)--(axis cs:3.4,-.85) node[midway,below,clearlabel]{$v_p$};
\end{axis}
\end{tikzpicture}
\figcap{窄带波包。载波相位在包络内移动，包络和相位的速度可以不同。}\label{ch12:fig:wave-packet}
\end{center}
图~\ref{ch12:fig:wave-packet} 中虚线包络的位置随 $v_g$ 移动，而某个波峰按 $v_p$ 穿过包络。

\begin{origexamplebox}{深水重力波的相速度与群速度}
\label{ch12:ex:06}
深水表面重力波满足 $kh\gg1$（$h$ 为水深），并处在重力远强于表面张力的波长范围。忽略黏性，色散关系 $\omega^2=gk$。求相速度和群速度的关系，并判断一个窄带波包中能量包络相对波峰怎样运动。
\end{origexamplebox}
\begin{xesolution}
取正频率支 $\omega=\sqrt{gk}$，
\[
 v_p=\frac\omega k=\sqrt{\frac gk},
 \qquad
 v_g=\frac{\dd\omega}{\dd k}=\frac12\sqrt{\frac gk}.
\]
所以
\[
 \boxed{v_g=\frac12v_p}.
\]
波峰以 $v_p$ 前进，会从波包后部生成、穿过包络并在前部消失；能量包络只以一半速度传播。
量纲上 $g/k$ 为 $L^2T^{-2}$，开方后是速度。若水深有限，色散关系改变，比例不再恒为 $1/2$。
\end{xesolution}

\subsection{从两个相近波数推导包络速度}
相速度跟踪一个确定相位，例如某一波峰。对 $kx-\omega t=\text{常量}$ 求时间导数，$k\dd x/\dd t-\omega=0$，所以 $v_p=\omega/k$。这并未回答一组有限长度波动的整体轮廓怎样移动。

先把问题简化为两个等振幅、同向的近频率行波：
\[
 y=A\cos(k_1x-\omega_1t)+A\cos(k_2x-\omega_2t).
\]
定义 $\bar k=(k_1+k_2)/2$、$\bar\omega=(\omega_1+\omega_2)/2$、$\Delta k=k_1-k_2$、$\Delta\omega=\omega_1-\omega_2$。和差化积给出
\[
 y=2A\cos\frac{\Delta kx-\Delta\omega t}{2}
       \cos(\bar kx-\bar\omega t).
\]
第二项快速振荡，第一项缓慢调节振幅。追踪一个包络极大，就要保持 $\Delta kx-\Delta\omega t$ 不变，于是
\[
 v_{\rm env}=\frac{\Delta\omega}{\Delta k}.
\]
当两波数都趋近中心波数 $k_0$，差商趋于色散曲线的斜率，才得到 $v_g=\omega'(k_0)$。因此相速度对应从原点到色散曲线上一点的割线斜率，群速度对应该点的切线斜率；只有线性无截距的 $\omega=ck$ 才在所有波数上都相同。

两个单色波的包络在整个空间周期重复，还不是真正局域的波包。若叠加一段窄波数范围，在 $k_0$ 附近展开
\[
 \omega(k)=\omega_0+v_g(k-k_0)+\tfrac12\omega''(k_0)(k-k_0)^2+\cdots,
\]
只保留常数项和一次项时，各分量的慢相位共同依赖 $x-v_gt$，局域包络近似刚性平移。二次项使不同波数获得不同额外相位，波包会展宽或变形；忽略它要求 $|\omega''(k_0)|(\Delta k)^2t\ll1$。图~\ref{ch12:fig:wave-packet} 中虚线轮廓对应这种窄带包络，内部细密波峰对应载波，不能把两个箭头解释为同一个质元的两种速度。

\begin{origexamplebox}{带截止频率的波包}\label{ch12:ex:teaching-cutoff}
某无耗一维线性系统的色散关系为 $\omega^2=\omega_c^2+c_0^2k^2$，其中 $\omega_c>0$、$c_0>0$ 均为常量。取沿 $+x$ 传播的窄带波包，中心频率为 $2\omega_c$。求相速度、群速度及二者之积，并讨论接近截止频率时的变化。
\end{origexamplebox}
\begin{xesolution}
取 $k>0$、$\omega>0$，在中心频率处有 $k_0=\sqrt3\omega_c/c_0$。先由相位定义计算
\[
 v_p=\frac{2\omega_c}{\sqrt3\omega_c/c_0}=\frac{2c_0}{\sqrt3}.
\]
对色散关系求导，$2\omega\,\dd\omega/\dd k=2c_0^2k$，因而
\[
 v_g=\frac{c_0^2k_0}{\omega_0}=\frac{\sqrt3c_0}{2},
 \qquad v_pv_g=c_0^2.
\]
中心频率从上方趋于 $\omega_c$ 时，$k_0\to0$，$v_p\to\infty$ 而 $v_g\to0$。发散的是等相位点的形式速度，不能据此断言信息以无限速度传递。对越来越靠近截止的波包，还须使带宽足够窄，才能继续使用单一群速度近似。
\end{xesolution}

\section{管中驻波}

开口端附近声压扰动近似为零，是声压节点、空气位移腹；刚性闭口端空气位移为零，是位移节点、声压腹。
两端开口或两端闭口的理想管长 $L$ 时
\begin{equation}
 f_n=\frac{nc}{2L},\qquad n=1,2,3,\ldots.
\end{equation}
一端闭口、一端开口时只允许奇次谐波：
\begin{equation}
 f_m=\frac{(2m-1)c}{4L},\qquad m=1,2,3,\ldots.
\end{equation}
真实开口的位移腹略在管口之外，可用端部修正增加有效长度。

\begin{center}
\begin{tikzpicture}[x=1.1cm,y=.8cm,>=Latex]
  \draw[very thick] (0,-1.2)--(0,1.2);
  \draw[thick] (0,1.2)--(6,1.2);
  \draw[thick] (0,-1.2)--(6,-1.2);
  \draw[DeepBlue,very thick,domain=0:6,samples=180] plot (\x,{.85*sin(45*\x)});
  \draw[guide] (0,0)--(6,0);
  \node[above] at (0,1.4) {闭口};
  \node[above] at (6,1.4) {开口};
  \draw[<->,TextGray] (0,-1.7)--(6,-1.7) node[midway,below]{$L=3\lambda/4$};
\end{tikzpicture}
\figcap{一端闭口一端开口管的第二个允许位移模态。曲线表示纵向位移随位置的分布，并非空气真的沿曲线弯曲；对应 $L=3\lambda/4$。}\label{ch12:fig:pipe-mode}
\end{center}
图~\ref{ch12:fig:pipe-mode} 的闭口是位移节点、开口是位移腹，因而管长只能容纳奇数个四分之一波长。

\begin{origexamplebox}{闭管的共振频率判别}
\label{ch12:ex:07}
一端闭口、一端开口的空气柱长 $L=0.85\,\mathrm m$，声速 $c=340\,\mathrm{m/s}$。求最低三个共振频率，并判断 $300\,\mathrm{Hz}$ 是否为共振频率。
\end{origexamplebox}
\begin{xesolution}
基频
\[
 f_1=\frac c{4L}=\frac{340}{4\times0.85}=100\,\mathrm{Hz}.
\]
闭管只取奇数倍，故最低三项为
\[
 \boxed{100\,\mathrm{Hz},\quad300\,\mathrm{Hz},\quad500\,\mathrm{Hz}}.
\]
因此 $300\,\mathrm{Hz}$ 是第二个允许模态，而不是“二次谐波”；按全部整数倍编号时它是第三谐波。
频率公式中 $c/L$ 的量纲为 $\mathrm{s^{-1}}$。
\end{xesolution}

\subsection{先决定边界约束的是位移还是声压}
令 $x$ 从闭口端指向开口端，空气纵向位移为 $\xi(x,t)$。闭口处空气不能穿过刚性壁，故 $\xi(0,t)=0$；开口处与外界相通，理想近似下声压扰动 $p'(L,t)=0$，并不是绝对压强为零。
对一个两端原本相距 $\dd x$ 的空气微元，实际长度增量为 $\xi_x\dd x$。质量守恒给出密度变化 $\delta\rho=-\rho_0\xi_x$，结合 $p'=c^2\delta\rho$，得到
\[
 p'=-\rho_0c^2\xi_x.
\]
因此开口的边界条件是 $\xi_x(L,t)=0$，与闭口的 $\xi=0$ 不同。对简谐驻波写 $\xi=X(x)\cos\omega t$，波动方程化为 $X''+k^2X=0$，通解为 $C\sin kx+D\cos kx$。闭口条件消去 $D$，开口条件给
\[
 X'(L)=Ck\cos kL=0,
 \quad kL=\frac{(2m-1)\pi}{2}.
\]
再用 $\omega=ck$，就得到闭管只允许奇次谐波的频率公式。图~\ref{ch12:fig:pipe-mode} 选的是 $m=2$：$kL=3\pi/2$，故 $L=3\lambda/4$。图中的曲线是“各位置的纵向位移值”，不是空气微元走过的弯曲路径。

同一驻波的声压空间因子正比于 $\cos kx$，而位移因子为 $\sin kx$。所以位移波节处为声压波腹，位移波腹处为声压波节；不能在同一位置把两者都标成波节。两端闭口给 $\sin kL=0$；两端开口可从 $X=C\cos kx$ 出发，也得到 $kL=n\pi$。零频率的整体平移不是这里讨论的声学共振模态。

\begin{origexamplebox}{由连续共振长度测声速与端部修正}\label{ch12:ex:teaching-resonance-length}
频率 $f=500\,\mathrm{Hz}$ 的音叉放在竖直细管口附近，管下部水面形成近似刚性闭端。调节水面，测得连续两个共振空气柱长度为 $L_1=0.160\,\mathrm m$、$L_2=0.500\,\mathrm m$。设两次开口端修正量相同，记为 $e$；且第一次为基模。求声速和 $e$。
\end{origexamplebox}
\begin{xesolution}
开口位移腹稍在管外，所以有效长度是 $L+e$。第一、第二允许模态分别满足
\[
 L_1+e=\frac\lambda4,\qquad L_2+e=\frac{3\lambda}4.
\]
两式相减可消去不能直接量取的 $e$：$L_2-L_1=\lambda/2$。因此
\[
 \lambda=0.680\,\mathrm m,\qquad
 c=f\lambda=340\,\mathrm{m/s},\qquad
 e=\lambda/4-L_1=0.010\,\mathrm m.
\]
修正量为正，与位移腹在管外一致。若不知道第一次是哪一级，仍能由相邻共振间隔求 $c$，却不能唯一确定 $e$；这说明“第一次为基模”不是多余条件。
\end{xesolution}

\section{介质边界的反射与透射}

两根张力相同而线密度不同的弦在 $x=0$ 连接。令特征阻抗 $Z_j=\sqrt{T\mu_j}$，位移振幅满足边界处位移连续和横向力连续。
用同一时间因子，写
$y_1=A_i e^{\ii(\omega t-k_1x)}+A_r e^{\ii(\omega t+k_1x)}$、
$y_2=A_t e^{\ii(\omega t-k_2x)}$，实际取实部。
$x=0$ 处位移连续要求 $A_i+A_r=A_t$，横向力连续要求
$k_1(A_i-A_r)=k_2A_t$。因 $Tk_j=\omega Z_j$，
第二式为 $Z_1(A_i-A_r)=Z_2A_t$。
将 $A_t=A_i+A_r$ 代入并移项，得到 $(Z_1-Z_2)A_i=(Z_1+Z_2)A_r$，
因此位移振幅系数为
\begin{equation}
 r=\frac{A_r}{A_i}=\frac{Z_1-Z_2}{Z_1+Z_2},
 \qquad
 t=\frac{A_t}{A_i}=\frac{2Z_1}{Z_1+Z_2}.
\end{equation}
功率系数为
\begin{equation}
 \mathcal R=r^2,
 \qquad
 \mathcal T=\frac{Z_2}{Z_1}t^2
 =\frac{4Z_1Z_2}{(Z_1+Z_2)^2},
 \qquad \mathcal R+\mathcal T=1.
\end{equation}
位移振幅可以大于入射振幅，仍不代表能量超过入射能量；功率还含阻抗因子。

\begin{center}
\begin{tikzpicture}[scale=.9,>=Latex]
  \draw[very thick,DeepBlue] (-4,0)--(0,0);
  \draw[very thick,TextGray] (0,0)--(4,0);
  \draw[vector] (-3,1.0)--(-.4,1.0) node[midway,above]{$A_i$};
  \draw[vector] (-.4,-.9)--(-3,-.9) node[midway,below]{$A_r$};
  \draw[vector] (.35,1.0)--(3,1.0) node[midway,above]{$A_t$};
  \draw[guide] (0,-1.4)--(0,2.2);
  \node[clearlabel] at (-2,2.15) {$Z_1$};
  \node[clearlabel] at (2,2.15) {$Z_2$};
  \node[point,label={[clearlabel]below right:连接点}] at (0,0) {};
\end{tikzpicture}
\figcap{弦连接点处的入射、反射与透射方向。连接点位移和横向力必须同时连续。}\label{ch12:fig:boundary}
\end{center}
按图~\ref{ch12:fig:boundary} 的传播方向，反射波的波数分量变号，这正是横向力边界式中出现 $A_i-A_r$ 的原因。

\begin{origexamplebox}{轻弦到重弦的反射相位与能量}
\label{ch12:ex:08}
同一张力下两弦线密度满足 $\mu_2=4\mu_1$，波从弦1入射到弦2。求位移振幅反射、透射系数和能量反射、透射率，并判断反射是否反相。
\end{origexamplebox}
\begin{xesolution}
$Z\propto\sqrt\mu$，所以 $Z_2=2Z_1$。于是
\[
 r=\frac{Z_1-Z_2}{Z_1+Z_2}=-\frac13,
 \qquad
 t=\frac{2Z_1}{Z_1+Z_2}=\frac23.
\]
$r<0$ 表示位移反射波反相。能量系数
\[
 \boxed{\mathcal R=r^2=\frac19},
 \qquad
 \boxed{\mathcal T=\frac{Z_2}{Z_1}t^2
 =2\times\frac49=\frac89}.
\]
二者之和为1，验证无耗边界的能量守恒。若 $Z_2\to\infty$，$r\to-1$，回到固定端反射；若 $Z_2=Z_1$，反射消失。
\end{xesolution}
